\documentclass[book.tex]{subfiles}

\begin{document}
\chapter{Quantum Advantage in Machine Learning}
\label{chap:quantum_advantage}
\noindent\rule{11cm}{0.4pt} \\
\textbf{Prerequisites:} \autoref{chap:molecular_hamiltonian}, \autoref{chap:dft}, \autoref{chap:hartree_fock} \\
\textbf{Difficulty Level:} ** \\
\noindent\rule{11cm}{0.4pt}
\vspace{5mm}

"The use of quantum computing for machine learning (ML) is among the most exciting prospective applications of quantum technologies. In this talk, I will present recent results for characterizing when quantum advantage in machine learning can be found. We will begin with theoretical results for understanding whether one can achieve good prediction performance by learning from significantly fewer quantum data compared to classical data. Then, we will restrict ourselves to ML problems with classical data and study the potential for computational advantages. We show that some problems that are classically hard to compute can be easily predicted by classical machines learning from data. Using rigorous prediction error bounds as a foundation, we develop a methodology for assessing potential quantum advantage in learning tasks. These constructions explain numerical results showing that with the help of data, classical machine learning models can be competitive with quantum models even on quantum many-body problems. We then propose improvements to existing quantum ML models that provide a simple and rigorous quantum speed-up for a learning problem in the fault-tolerant regime. For near-term implementations, we demonstrate a significant prediction advantage over various classical ML models on engineered data sets in one of the largest numerical tests for gate-based quantum machine learning to date, up to 30 qubits."

References \cite{huang2021information}, \cite{huang2021power}

\end{document}